\documentclass[10pt,a4paper]{amsart}
\usepackage{amsmath,amssymb,amsthm,graphicx}
\usepackage{amsmath}
\usepackage{amssymb}
\usepackage{hyperref}
\renewcommand{\theenumi}{\Alph{enumi}}
\renewcommand{\thesubsection}{\thesection\Alph{subsection}}

%%%%%%%%%%%%%%%%%%%%%%%%%%%%%%%%%%%%%%%%%%
% ENVIRONMENTS
%%%%%%%%%%%%%%%%%%%%%%%%%%%%%%%%%%%%%%%%%%%
\newtheorem{theorem}{Theorem}[section]
\newtheorem{prop}[theorem]{Proposition}
\newtheorem{lemma}[theorem]{Lemma}
\newtheorem{remark}[theorem]{Remark}
\newtheorem{corollary}[theorem]{Corollary}
\newtheorem{definition}[theorem]{Definition}
\newtheorem{example}[theorem]{Example}
%
\def\beq{\begin{equation}}
\def\eeq{\end{equation}}
\def\be{\begin{equation*}}
\def\ee{\end{equation*}}
%%%%%%%%%%%%%%%%%%%%%%%%%
% Beginning of the text %
%%%%%%%%%%%%%%%%%%%%%%%%%
%\begin{document}



%\documentclass[10pt]{article}
\usepackage{amssymb}
\usepackage{amsfonts}
%\topmargin 0.3in
%\textwidth 6in
%\textheight 9in%\baselineskip
%\hoffset=-50pt
%\voffset=-2pt
%\newenvironment{proof}[1][Proof]
%{\ \rule{0.5em}{0.5em}}


\title{Dedekind-finite Cardinals having countable partitions}
\author{Supakun Panasawatwong and J K Truss} 
\date{This paper is an amplified version of part of the first author's PhD thesis at the University of Leeds}%\footnotemark}
\title{Every template encodes a finitely branching omega-tree, uniquely up to local isomorphism, with branches controlled by branches of the template}
\author{Supakun Panasawatwong, J. K. Truss}
\date{arxiv:2204.05700v1; CC BY 4.0}
\begin{document}
\maketitle
\noindent\textit{Editorial note.} All delivered mathematical text is the author's own, copied from the pinned publisher source sources/dedekind-finite-trees-2204.05700v1/source0.tex; no statement or proof step was rewritten, reordered or removed.\par\medskip
\noindent\textit{Editorial note.} Preamble provenance: the retained preamble is the source's own layout adapter, if any, plus the author's macro block at lines 5-44 of sources/dedekind-finite-trees-2204.05700v1/source0.tex (everything between the publisher class metadata and \textbackslash{}begin\{document\}, comments excluded). Nothing was added, no macro was invented, and the publisher document class is deliberately not reproduced; the excerpt serves only as preamble, never as a span coordinate.\par\medskip
\noindent\textit{Editorial note.} External references are navigation pointers only; every deliberately kept step of the author's argument stays inside the counted spans. The displayed numbers of those pointers are the author's own numbering under the author's \textbackslash{}newtheorem declarations (the source has no subsections, so numbering is section.index), computed mechanically by the lane's scanner.\par\medskip
\noindent\textit{Editorial note.} The counted proof contains no reference to any labelled item and no bibliography citation, so the card needs neither a navigation stub nor a bibliography excerpt.\par\medskip
\noindent\textit{Editorial note.} Source metadata: this directory stages no SOURCE.json, so the title and the author list were taken from the publisher abstract page kept beside it, and the identifiers from the staged licence evidence; the card points at a per-card licence reuse record in its own task directory because that evidence carries no status field.\par\medskip
\noindent\textit{Allowed standard prerequisites (editorial).}\begin{itemize}
\item The notion of a template and of a tree encoded by it, with the levels P\_n and the attached positive integers. Set up in the definitions preceding the counted statement (the notation P\_n, L\_n and the partition correspondence used in the proof). Fixes the data that the four clauses quantify over and the meaning of local embeddability used in clause (ii).
\item The Fraenkel-Mostowski construction with finite supports used earlier in the paper. Used inside the counted proof of clause (iv) at lines 620-625. Supplies the permutation model in which a tree encoded by a template without an eventually singleton branch has no infinite branch.
\end{itemize}\medskip
\begin{lemma} (i) For any template $(T^*, \prec)$ there is a finitely branching $\omega$-tree $(T, <)$ encoded by $(T^*, \prec)$.

(ii) Any two trees encoded by $(T^*, \prec)$ are locally isomorphic.

(iii) If $(T^*, \prec)$ has an eventually singleton $\omega$-branch, then $(T, <)$ has an infinite branch.

(iv) If $(T^*, \prec)$ is a template having no eventually singleton branch, then in a suitable Fraenkel--Mostowski model, 
there is a tree encoded by $(T^*, \prec)$ having no infinite branch. \label{4.3}  \end{lemma}

\begin{proof} (i) Let $(T^*, \prec)$ be a template with $n$th level $P_n$ and positive integers attached as in the
definition. We build $(T,<)$ level by level, using induction. There will be a bijection between $P_n$ and a partition $\pi_n$ of
$L_n$. Choose any root for $T$, and let $\pi_0$ be the trivial singleton partition of $L_0$ which just contains the root, and
$P_0$ corresponds to $\pi_0$.

Assuming that $L_n$ has been chosen, as has a partition $\pi_n$ of $L_n$ which corresponds to $P_n$ under some bijection, for
$X \in P_n$ let $Y_1, \ldots, Y_k$ be the members of $P_{n+1}$ above $X$ under $\prec$, with positive integers $n_1, \ldots, n_k$
attached. Under the bijection given by the induction hypothesis, $X$ may be viewed as a subset of $L_n$. For each $x \in X$, and
$i$ between 1 and $k$, choose $n_i$ new points which together (as $i$ varies) will form $succ(x)$ in $T$. Doing this for all
$x$ determines $L_{n+1}$. The $(n+1)$th level $P_{n+1}$ of $T^*$ is now identified with a partition $\pi_{n+1}$ of $L_{n+1}$ by placing
$y_1$ and $y_2$ in the same piece of $\pi_{n+1}$ provided that for some $X \in \pi_n$, $x_1, x_2 \in X$, and $i$ between 1 and $k$,
$y_1$ and $y_2$ are adjoined corresponding to that value of $i$.

(ii) We observe that we have the freedom to choose the new points adjoined with or without an ordering. Thus if we are (secretly) 
thinking of them as ordered, then the whole of the tree can be well-ordered, and has cardinality $\aleph_0$. If we make no attempt
to order them, then there is no reason why this should be true. So we do not expect the two possibilities to give rise to 
isomorphic trees. It is clear however that they will be locally isomorphic, since their behaviour is identical up to any finite
stage.

(iii) As remarked above, a branch of the template corresponds to a balanced subtree of $T$, and saying that this is eventually singleton
just says that this subtree is eventually just a single branch.

(iv) What we really want to say is that $(T,<)$ has an infinite branch if and only if its template has an eventually singleton branch,
but for the reasons explained this is too strong. Suppose then that $(T^*, \prec)$ is a template and that it encodes $(T, <)$. We form
a Fraenkel--Mostowski model in which $U$ is indexed by the members of $T$, and we let $G$ be the group of automorphisms of $U$ induced by
tree automorphisms of $(T,<)$ which fix each member of $T^*$, and we let $\mathfrak N$ be the FM-model defined using finite supports.
One sees that $G$ acts on each subtree obtained as above from an $\omega$-branch of $T^*$. This is a balanced tree which has
infinitely many non-singleton entries, and so as for the simplest construction given from just one balanced tree, has no infinite
branch in $\mathfrak N$. Since any infinite branch of $T$ would have to arise in this way, there cannot be any at all. \end{proof}

\end{document}
